\documentclass[twocolumn]{iopjournal}

\usepackage{amsmath,amssymb,amsfonts}
\usepackage{amsthm}

\fancyhfoffset[L]{0mm}
\setlength{\marginparwidth}{0pt}
\setlength{\marginparsep}{0pt}

% ==========================================
% ZENODO PREPRINT HEADERS (ACTIVE)
% ==========================================
\renewcommand{\articletype}[1]{%
  {\vspace*{-8mm}\noindent \Large \sf Preprint}\par
  \vspace*{6mm}\noindent{\scriptsize \sf{\bfseries \MakeUppercase{#1}}}}

\fancyhead[L]{}
\fancyhead[R]{}

\hypersetup{pdftitle={The Net-Field Rule: Where the Constitutive Law Acts}, pdfauthor={Stephen Nelson}, pdfsubject={Gravity Field Dynamics (GFD)}, pdfkeywords={constitutive law, net-field rule, external field effect, MOND, solar system, Cassini, globular clusters, wide binaries, proof}}

\newtheorem{lemma}{Lemma}
\newtheorem{theorem}{Theorem}
\newtheorem{proposition}{Proposition}
\newtheorem{corollary}{Corollary}
\theoremstyle{definition}
\newtheorem{definition}{Definition}
\theoremstyle{remark}
\newtheorem*{remark}{Remark}

\begin{document}

\articletype{Proof}

\title{The Net-Field Rule: Where the Constitutive Law Acts}

\author{Stephen Nelson}

\affil{Field Dynamics Research}

\ead{stephen@fielddynamics.org}

\begin{abstract}
The constitutive law of Gravity Field Dynamics, $\mu(x) = x/(x+1)$,
fits galaxies without dark matter, but applied to each star's own field
it leaves a pull of nearly $a_0$ at every planet and an external-field
quadrupole five times the Cassini bound. We prove that this failure is
forced for every exchange-symmetric law with flat rotation curves. The
rule that avoids it rests on two premises, which we state rather than
derive. First, a drain belongs to a field, not to a body: $\mu$ acts only
on a field that has a drain, a galaxy's, while every body's TetB is
enclosed in it and counted in its $GM$. Second, $\mu$ acts on that
field's net field coarse-grained on the tidal (Jacobi) scale of each
bound system. Given these, a field without a drain is Newton and general
relativity with the bare $G$, and four properties are equivalent: the
law is linear, a body inside responds with the bare $G$, flux is
conserved across empty shells, and there is no TetB share. Each is a
rational function of $\mu$ vanishing only at $\mu = 1$, and the flux
exponent $2(1-\mu)/(2-\mu)$ is an involution of $\mu$ that singles out
$x/(x+1)$ among the standard laws. A star system is a rubber sheet in its
galaxy's field: inside it the galaxy acts only through free fall and its
tide, about $10^{-18}$\,m\,s$^{-2}$ at Saturn, and Cassini puts the
sheet's edge beyond $1.2 \times 10^4$\,AU, consistent with the Jacobi
radius of $2.9 \times 10^5$\,AU. Two remote globular clusters require
the bare $G$ at more than $8\sigma$, and bound wide binaries are
predicted to be Newtonian to within $(s/r_J)^3$.
\end{abstract}

\keywords{constitutive law, net-field rule, external field effect, MOND,
solar system, Cassini, globular clusters, wide binaries, proof}

% ======================================================================

\section*{Setting}

In a galaxy, the total pull $g_t$ and the Newtonian pull of the visible
matter $g_N$ are related by $\mu(g_t/a_0)\,g_t = g_N$, with
$\mu(x) = x/(x+1)$ the share of the field carried by TetA and $1 - \mu$
the share carried by TetB~\cite{NelsonL2,NelsonL3,NelsonTug}. Solving,
\begin{gather}
  g_t = F(g_N) = \tfrac12\Bigl(g_N + \sqrt{g_N^2 + 4 g_N a_0}\Bigr),
  \label{eq:F}\\
  g_N = \frac{g_t^2}{g_t + a_0} = \mu(g_t/a_0)\,g_t. \nonumber
\end{gather}
The law fits the rotation curves of galaxies with one measured
acceleration, $a_0$~\cite{Milgrom1983,SPARC}. The question this proof settles is
\emph{which} field it acts on. AQUAL~\cite{BM1984} and
QUMOND~\cite{Milgrom2010} apply it to the total field at every point,
and that reading is excluded in the solar
system~\cite{Hees2014,Hees2016}. The framework's answer, which we make
exact, rests on three premises.

\begin{definition}[premises]
\textbf{(D) Drain criterion.} A drain belongs to a field, not to a
body. Every body has its own TetB, enclosed within it and counted in its
$GM$. A field carries a TetB share, and $\mu$ acts on it, only if the
field itself has a drain: a throat opening into occupied TetB. A
galaxy's field has one. The field of a star system, a binary or a star
cluster does not.
\textbf{(N) Net field.} In a drained field, $x = |g_{\rm net}|/a_0$, with
$g_{\rm net}$ that field's own net field, coarse-grained on the tidal
(Jacobi) scale $\ell = r_J$ of each bound system inside it.
\textbf{(A) TetA is general relativity.} A drainless field is the TetA
connection: Newton and GR, with the bare $G$.
\end{definition}

The picture is a rubber sheet with a bowling ball on it, sitting in the
galaxy's stream. The ball, the Sun, has a TetB, enclosed in it. The
sheet, the solar system's field, has no drain. A galaxy's field is a
funnel whose throat opens into TetB. Premises (D) and (N) are the
framework's~\cite{NelsonEinstein,NelsonBimetric}, and the scale in (N)
is motivated in Proposition~\ref{prop:edge} but not derived; what
follows proves what the premises imply and tests it against the data
the framework must meet.

\section*{Four faces of a drainless field}

\begin{theorem}[the four faces]\label{thm:faces}
Let $m = \mu(x) \in (0,1]$. For the law~\eqref{eq:F}:
\begin{enumerate}
\item[(i)] \emph{Linearity.} $F$ is strictly subadditive,
  $F(b + c) < F(b) + F(c)$, and its degree is
  $d\ln F/d\ln g_N = 1/(2 - m)$: one only at $m = 1$, one half in the deep
  field.
\item[(ii)] \emph{Response.} A guest field $\delta$ sent through the
  host's law at host field $g_e$ responds with $\nu_\perp = 1/m$
  across it and $\nu_\parallel = 1/[m(2-m)]$ along it. Both exceed 1 at
  every finite field.
\item[(iii)] \emph{Flux.} Around a point mass the flux
  $\Phi = 4\pi r^2 g_t$ grows as $r^s$ with
  \begin{equation}
    s = \frac{2(1-m)}{2-m}, \qquad
    \frac{d\ln v^2}{d\ln r} = -\frac{m}{2-m}. \label{eq:s}
  \end{equation}
  The excess flux is $M(1-m)/m = M/\Omega$, the mass called dark matter.
  For any configuration of sources a drained field has nonzero divergence
  in empty space, and a drainless one has none.
\item[(iv)] \emph{TetB share.} $1 - m$.
\end{enumerate}
The four defects $1 - 1/(2-m)$, $1/m - 1$, $s$ and $1 - m$ vanish
together, exactly at $m = 1$.
\end{theorem}

\begin{proof}
(i) Squaring, subadditivity reduces to $bc < \sqrt{b^2+4ba_0}\,
\sqrt{c^2+4ca_0}$, true since each root exceeds its first term. The
degree is the reciprocal of $d\ln g_N/d\ln g_t = 2 - g_t/(g_t + a_0) =
2 - m$. (ii) Linearising $g = \nu(|g_N|/a_0)\,g_N$ about $g_e\hat e$
gives $\nu\,\delta_\perp + (\nu + y\nu')\,\delta_\parallel$. With
$y = x\mu(x)$, $1/4 + 1/y = [(x+2)/2x]^2$, so $\nu = (x+1)/x = 1/m$, and
$\nu + y\nu' = dg_t/dg_N = (x+1)^2/[x(x+2)] = 1/[m(2-m)]$. (iii) With
$g_N = GM/r^2$, $d\ln g_t/d\ln r = -2/(2-m)$, so $s = 2 - 2/(2-m)$. The
excess is $\Phi/4\pi G - M = M/m - M$. A sum of point fields is
divergence-free off the sources; $\nu g_N$ is not (checked at random
points). The last statement is immediate.
\end{proof}

\begin{corollary}[the flux involution]\label{cor:inv}
The map $f(u) = 2(1-u)/(2-u)$ that sends $m$ to $s$ satisfies
$f(f(u)) = u$. Outside the baryons, where the enclosed mass is constant,
the slope of a rotation curve alone gives $m = f(s)$ with
$s = 1 + d\ln v^2/d\ln r$, then $x = m/(1-m)$ and $a_0 = (v^2/r)/x$. In
the standard family $\mu_n = x/(1+x^n)^{1/n}$, where
$x\mu_n'/\mu_n = 1 - \mu_n^n$ and $s = 2(1-\mu^n)/(2-\mu^n)$, the map is
an involution only for $n = 1$.
\end{corollary}

\begin{proof}
$1 - f(u) = u/(2-u)$ and $2 - f(u) = 2/(2-u)$, so $f(f(u)) = u$. For
$n = 2$, $f_2(f_2(0.3)) \neq 0.3$. On a synthetic galaxy the slope
returns $a_0$ to $0.2\%$.
\end{proof}

\section*{A body's TetB}

\begin{proposition}[enclosed, screened, bounded]\label{prop:body}
(a) Outside a body, the field depends only on its total mass, however
that mass is split between TetA and an enclosed TetB, so no orbit can
separate them. Flows inside it, such as a sunspot's vortex closing
within the Sun, conserve the enclosed mass and change the exterior only
through multipoles $\propto r^{-(l+1)}$. (b) Under the screening law
$x(1-\mu) = \mu$, the TetB share is $1/(1+x)$: about
$4 \times 10^{-13}$ at the Sun's surface, rising to 1 only in a core of
radius $r_B = 3a_0/(4\pi G\rho_c) \approx 3\,\mu$m at the centre, the
ball's Field Origin. (c) Read pointwise, the screened TetB's extra mass
is below $a_0R_\odot^2/G \approx 9 \times 10^{17}$\,kg, or
$4 \times 10^{-13} M_\odot$, fixed by $R_\odot$ and counted in $GM$. The
ephemerides give $\dot{(GM)}/GM = (-5.0 \pm 4.1)\times10^{-14}$ per year,
which radiation alone ($-6.8\times10^{-14}$) accounts
for~\cite{Pitjeva2012}, so any TetB exchange across the surface is below
about $10^{-13}$ per year.
\end{proposition}

\begin{proof}
(a) Gauss's law, and continuity with $v\cdot n = 0$ at the surface;
each multipole solves Laplace's equation. (b) $1 - x/(x+1) = 1/(1+x)$;
in a uniform core $g = \tfrac43\pi G\rho_c r$. (c) By
Proposition~\ref{prop:tail}, $g_t - g_N < a_0$ at the surface.
\end{proof}

\section*{A drain around every star fails}

\begin{proposition}[forced failure]\label{prop:tail}
(a) Under~\eqref{eq:F}, a star's field gains a pull with
$a_0\,g_N/(g_N + a_0) \le g_t - g_N < a_0$: nearly $a_0$ at every planet,
a constant pull that no rescaling of $GM$ removes. (b) This is forced.
For every $\mu$ with the exchange $\mu(1/x) = 1 - \mu(x)$ and the deep
limit $\mu(u)/u \to 1$ as $u \to 0^+$, $x(1-\mu(x)) = \mu(1/x)/(1/x)
\to 1$ as $x \to \infty$. (c) In the Milky Way's field at the Sun,
$x \approx 1.9$, the law applied pointwise gives an external-field
quadrupole $Q_2 \approx 4.6 \times 10^{-26}$\,s$^{-2}$, against
Cassini's $(3 \pm 3) \times 10^{-27}$\,s$^{-2}$~\cite{Hees2014}.
\end{proposition}

\begin{proof}
(a) $g_t - g_N = 2g_Na_0/(\sqrt{g_N^2 + 4g_Na_0} + g_N)$, and
$g_N < \sqrt{\cdot} < g_N + 2a_0$. (b) Substitute $u = 1/x$ in the
exchange. (c) The phantom density of the QUMOND form, integrated for
the tidal component at the Sun~\cite{Milgrom2010,NelsonBimetric}.
\end{proof}

\section*{The rubber sheet}

Let $P$ be the pointer projection: replace the mass inside each sheet, a
ball of radius $\ell$ about its centre of mass, by a uniform ball of the
same mass. Write the field as
\begin{equation}
  g = F\bigl[g_N(P\rho)\bigr] + g_N(\rho - P\rho). \label{eq:split}
\end{equation}
By linearity of $g_N$ this is an identity; the rule is that only the
first term, the drained host, goes through $F$, while the second, the
sheet's own field, stays linear with the bare $G$ by (A).

\begin{theorem}[the rubber sheet]\label{thm:sheet}
(a) $P$ conserves mass, is idempotent, and nested projections compose:
projecting a small ball and then a larger one containing it is the same
as projecting the larger one. Far away, a sheet's field is its monopole
at the centre of mass, with no dipole and corrections of order
$(d/D)^2$. (b) Outside the sheet, $\rho - P\rho$ has no monopole and, about
the centre of mass, no dipole. For a spherical sheet its field vanishes
identically, $GM/r^2 - g_{\rm ball}(r) = 0$ for $r \ge \ell$, and is
continuous at $\ell$; inside, it is $GM/r^2 - GMr/\ell^3$. Otherwise what
remains outside is the sheet's own quadrupole and higher multipoles,
falling as $r^{-4}$ or faster, ordinary Newtonian gravity counted once
(for the solar system at 1\,pc, $6\times10^{-13}$ of the monopole). The
host's field, and every galactic prediction, is unchanged.
(c) A uniform host field drops out of every relative acceleration, and
with pairwise antisymmetric internal forces the centre of mass falls
with it. (d) No self-force: the term $k(p - c)$, $k = GM/\ell^3$, that
the host term returns inside the sheet exerts zero net force for any
linear host response $R$, $\sum_i m_i R\,k(p_i - c) = 0$. What remains
inside is $(R - I)\,k\,(p - c)$ plus the tide $-\nabla\nabla\Phi\cdot d$.
(e) For a flat rotation curve, $\Phi = v^2\ln R$, the tide has size
$v^2/R^2 \approx 8 \times 10^{-31}$\,s$^{-2}$ at the Sun, about
$1.2 \times 10^{-18}$\,m\,s$^{-2}$ at Saturn, four orders below Cassini,
and set by the rotation curve alone, so every model of the Milky Way
shares it. The epicyclic frequency is $\kappa^2 = 2\Omega^2$, and at the
Jacobi radius $r_J^3 = GM/2\Omega^2$ the residual's scale is exactly the
tide's, $k = 2\Omega^2$.
\end{theorem}

\begin{proof}
(a) $P$ depends only on the mass inside its ball; expand the potential
of the sheet's masses about the centre of mass. (b) The shell theorem,
and $\sum_i m_i(p_i - c) = 0$.
(c) The host term cancels in differences; $\sum_{ij}F_{ij} = 0$ for
$F_{ij} = -F_{ji}$. (d) $\sum_i m_i(p_i - c) = 0$ and $R$ is linear. A
Taylor expansion gives the tide. (e) $\Phi'' = -v^2/R^2$,
$\Phi'/R = v^2/R^2$, and $\kappa^2 = R\,d\Omega^2/dR + 4\Omega^2$.
\end{proof}

\begin{proposition}[the sheet's edge]\label{prop:edge}
The residual's anisotropic part, $|\nu_\perp - \nu_\parallel|\,GM/\ell^3$,
is a quadrupole, $0.39\,GM/\ell^3$ at the Sun. Cassini requires
$\ell \gtrsim 1.2 \times 10^4$\,AU, about $1.7\,r_M$ with
$r_M = \sqrt{GM/a_0} \approx 7000$\,AU; a sheet ending at the Sun's own
ring $x = 1$ leaves $4.5 \times 10^{-26}$\,s$^{-2}$, the pointwise reading
again. Binding caps the edge at the Jacobi radius, $1.4$\,pc or
$2.9 \times 10^5$\,AU, 24 times the bound. Both scale as $M^{1/3}$, so the
window is the same for every star at the Sun's radius. From $GM$, $\Omega$
and $a_0$ every length is $r_M f(\pi)$ with the one group
$\pi = GM\Omega^4/a_0^3$; of the simple choices, $r_M$ is excluded,
$a_0/\Omega^2 \approx 5$\,kpc lies beyond binding, and only $r_J$, the
tidal edge, falls in the window. Whatever the edge's shape, a kernel
whose density never exceeds $c$ times the uniform ball's leaves a
residual at most $c(\nu-1)(r/\ell)^3$ of the system's own pull, so
predictions for $r \ll \ell$ do not depend on it.
\end{proposition}

\section*{The data}

\begin{proposition}[bare $G$, and the drain criterion]\label{prop:data}
(a) A host factor on $G$, $G \to G/\bar\mu$, is invisible to orbits,
since Kepler's law sees only $GM$. (b) Photometric masses see it. The
outer-halo clusters Pal~14 and Pal~4~\cite{Jordi2009,Frank2012} sit at
$x \approx 0.13$ and $0.12$. With the bare $G$ their dispersions fit
within $1.0\sigma$ and $0.4\sigma$; with the host factor
($\nu \approx 8$) they are excluded at $8.9\sigma$ and $11.1\sigma$. So
the host enters a drainless field only through free fall and tide, never
as a factor on $G$. (c) The dwarf Crater~II~\cite{Caldwell2017,Ji2021},
at the same $x \approx 0.08$ in the same host, is $5.4\sigma$ above
Newton and $2.2\sigma$ from $\mu$ with the host's field. No rule that
depends on acceleration alone fits all three ($\chi^2 = 29.9$ for the
best common factor); the drain criterion, clusters closed and the dwarf
open, fits with $\chi^2 = 6.1$ and no parameter.
\end{proposition}

\begin{proposition}[predictions]\label{prop:pred}
(a) With $\ell = r_J$, a bound wide binary is Newtonian to within
$0.53(s/\ell)^3$, below $10^{-4}$ at 20\,kAU, including beyond
$r_M \approx 10^4$\,AU. A boost growing as $s^3$ inside a cutoff would
instead measure $\ell$: $15\%$ at 10\,kAU if $\ell$ sat at the Cassini
bound. Rules that send the binary through the host's
law give $1.17$, $1.35$ and $1.47$ at $5$, $10$ and $20$\,kAU for
$2\,M_\odot$, the boost claimed in~\cite{Chae2023} and disputed
in~\cite{Banik2024}. (b) The local slope of the radial acceleration
relation is the inverse degree of the law, $\beta = 1/(2-\mu)$, with
$\mu = g_{\rm bar}/g_{\rm obs}$ measured: a test of the law's shape that
uses no value of $a_0$. On 3284 SPARC points in 163 galaxies, binned in
$g_{\rm bar}$, GFD's shape misses the measured slope by $0.053$ rms and
the standard law's $1/(2-\mu^2)$ by $0.102$, and GFD's wins in 200 of
200 bootstrap resamplings of the galaxies.
\end{proposition}

\begin{remark}[the vortex]
The picture is exact fluid mechanics. Navier--Stokes is Galilean
invariant, so a pebble in a stream sees the stream's gradient, never the
stream: free fall and tide. An open drain sustains the steady Burgers
vortex~\cite{Burgers1948}; with no drain the Lamb--Oseen vortex decays,
keeping its total circulation, as a body keeps its $GM$.
\end{remark}

\section*{Scope}

The premises (D) and (N) are the framework's, including the
coarse-graining scale $\ell = r_J$, which Proposition~\ref{prop:edge}
motivates by elimination but does not derive. What the proof adds is
their consequences, made exact, and the data they must meet. Three
further points remain open: an observable marker of a drained field that does
not use its rotation curve (Theorem~\ref{thm:faces}(iii) says it must
measure flux growth across empty shells); the covariant derivation of
which time leg's acceleration enters $\mu$, left open
in~\cite{NelsonSST}; and the dwarf Antlia~2, which no rule tested here
fits, possibly because it is being tidally disrupted. The data tests use
simple dispersion estimators~\cite{Wolf2010} and binned medians; the
RAR slope is sensitive to scatter and to stellar mass-to-light ratios.

\section*{Verification}

The script \texttt{verify\_net\_field\_rule.py} checks every statement,
128 checks in all: exactly with sympy where it is algebraic, by
quadrature for the Cassini quadrupole, and against the data for
Propositions~\ref{prop:data} and~\ref{prop:pred}. Its output is in
\texttt{verify\_output.txt}.

% Lean: new paragraph.
\textbf{Formal verification.} Theorem~\ref{thm:faces} (the dictionary,
the four faces, subadditivity, the responses and the flux exponent),
Corollary~\ref{cor:inv} (the involution and its failure for $n = 2$),
Propositions~\ref{prop:body}(b)--(c) and~\ref{prop:tail}(a)--(b), and
Theorem~\ref{thm:sheet}(b)--(e), the kernel bound of
Proposition~\ref{prop:edge} and Proposition~\ref{prop:data}(a) are checked in Lean~4 with Mathlib
(\texttt{lean/NetFieldRule.lean}), importing the constitutive law and the
exchange of the tug-of-war proof. Every theorem uses only Lean's three
standard axioms. The quadrupole integral and the data are checked by the
script.

\begin{thebibliography}{99}

\bibitem{Milgrom1983}
M.~Milgrom, A modification of the Newtonian dynamics as a possible
alternative to the hidden mass hypothesis, \textit{Astrophys. J.}
\textbf{270}, 365 (1983).

\bibitem{BM1984}
J.~Bekenstein and M.~Milgrom, Does the missing mass problem signal the
breakdown of Newtonian gravity?, \textit{Astrophys. J.} \textbf{286}, 7
(1984).

\bibitem{Milgrom2010}
M.~Milgrom, Quasi-linear formulation of MOND, \textit{Mon. Not. R.
Astron. Soc.} \textbf{403}, 886 (2010).

\bibitem{Hees2014}
A.~Hees, W.~M.~Folkner, R.~A.~Jacobson and R.~S.~Park, Constraints on
modified Newtonian dynamics theories from radio tracking data of the
Cassini spacecraft, \textit{Phys. Rev. D} \textbf{89}, 102002 (2014).

\bibitem{Hees2016}
A.~Hees, B.~Famaey, G.~W.~Angus and G.~Gentile, Combined Solar System
and rotation curve constraints on MOND, \textit{Mon. Not. R. Astron.
Soc.} \textbf{455}, 449 (2016).

\bibitem{Pitjeva2012}
E.~V.~Pitjeva and N.~P.~Pitjev, Changes in the Sun's mass and gravitational
constant estimated using modern observations of planets and spacecraft,
\textit{Sol. Syst. Res.} \textbf{46}, 78 (2012).

\bibitem{Jordi2009}
K.~Jordi \textit{et al.}, Testing fundamental physics with distant star
clusters: analysis of observational data on Palomar~14,
\textit{Astron. J.} \textbf{137}, 4586 (2009).

\bibitem{Frank2012}
M.~J.~Frank \textit{et al.}, The velocity dispersion and mass function of
the outer halo globular cluster Palomar~4, \textit{Mon. Not. R. Astron.
Soc.} \textbf{423}, 2917 (2012).

\bibitem{Caldwell2017}
N.~Caldwell \textit{et al.}, Crater~2: an extremely cold dark matter
halo, \textit{Astrophys. J.} \textbf{839}, 20 (2017).

\bibitem{Ji2021}
A.~P.~Ji \textit{et al.}, Kinematics of Antlia~2 and Crater~2 from the
Southern Stellar Stream Spectroscopic Survey, \textit{Astrophys. J.}
\textbf{921}, 32 (2021).

\bibitem{Wolf2010}
J.~Wolf \textit{et al.}, Accurate masses for dispersion-supported
galaxies, \textit{Mon. Not. R. Astron. Soc.} \textbf{406}, 1220 (2010).

\bibitem{Chae2023}
K.-H.~Chae, Breakdown of the Newton--Einstein standard gravity at low
acceleration in internal dynamics of wide binary stars,
\textit{Astrophys. J.} \textbf{952}, 128 (2023).

\bibitem{Banik2024}
I.~Banik \textit{et al.}, Strong constraints on the gravitational law
from Gaia DR3 wide binaries, \textit{Mon. Not. R. Astron. Soc.}
\textbf{527}, 4573 (2024).

\bibitem{SPARC}
F.~Lelli, S.~S.~McGaugh and J.~M.~Schombert,
SPARC: mass models for 175 disk galaxies with Spitzer photometry at
3.6\,$\mu$m and accurate rotation curves, \textit{Astron. J.}
\textbf{152}, 157 (2016).

\bibitem{Burgers1948}
J.~M.~Burgers,
A mathematical model illustrating the theory of turbulence,
\textit{Adv. Appl. Mech.} \textbf{1}, 171 (1948).

\bibitem{NelsonL2}
S.~Nelson, \textit{The Other Half of Gravity: Torsion and the Two-Chamber
Structure of the Tetrad Frame} (Letter~2), Field Dynamics Research (2026).

\bibitem{NelsonL3}
S.~Nelson, \textit{Gravity Field Dynamics: Field Closure of the Tetrad
Frame} (Letter~3), Field Dynamics Research (2026).

\bibitem{NelsonTug}
S.~Nelson, \textit{The Constitutive Law as a Tug of War}, Field Dynamics
Research (2026), \url{proofs/constitutive-law-tug-of-war}.

\bibitem{NelsonSST}
S.~Nelson, \textit{The SST Covariant Completion: $y$ From the
Acceleration of the Time Leg}, Field Dynamics Research (2026),
\url{proofs/sst-covariant-time-leg}.

\bibitem{NelsonEinstein}
S.~Nelson, \textit{Einstein's Equation as the TetA Partition}, Field
Dynamics Research (2026), \url{EINSTEIN_TETA_PARTITION.md}.

\bibitem{NelsonBimetric}
S.~Nelson, \textit{Does the Locked Bimetric Potential Project onto
$\mu(x) = x/(x+1)$?}, working note, Part~3, Field Dynamics Research
(2026), \url{research/bimetric-projection}.

\end{thebibliography}

\end{document}
